\Needspace{9\baselineskip}\begin{center}\itshape \textit{Proof of the theorem mentioned in Section IV.}\end{center}


\Needspace{6\baselineskip}\Article{357}
Let the equation of degree \(m\), whose roots are the \(m\) roots contained in the period \((m,1)\), be \[x^m-ax^{m-1}+bx^{m-2}-\text{etc.}=0\] or \(z=0\). Then \(a=(m,1)\), and each of the remaining coefficients \(b\), etc., is contained in a form such as \(\mathfrak A+\mathfrak B(m,1)+\mathfrak C(m,g)\), with \(\mathfrak A,\allowbreak \mathfrak B,\allowbreak \mathfrak C\) integers (Article 348). Denoting by \(z'\) the function into which \(z\) passes when \((m,1)\) is everywhere replaced by \((m,g)\), and \((m,g)\) by \((m,gg)\), or equivalently \((m,1)\), the roots of the equation \(z'=0\) will be the roots contained in \((m,g)\), and the product is \[zz'=\frac{x^n-1}{x-1}=X\]

Thus \(z\) can be reduced to a form such as \(R+S(m,1)+T(m,g)\), where \(R,\allowbreak S,\allowbreak T\) will be polynomials in \(x\), all of whose coefficients will also be integers. This gives \[z'=R+S(m,g)+T(m,1)\]

From this, writing for brevity \(p\) and \(q\) for \((m,1)\) and \((m,g)\), respectively, we obtain \[\begin{aligned}2z&=2R+(S+T)(p+q)-(T-S)(p-q)\\&=2R-S-T-(T-S)(p-q)\end{aligned}\] and similarly \[2z'=2R-S-T+(T-S)(p-q)\] Thus, putting \[2R-S-T=Y,\quad T-S=Z\]
\SourceLeaf{449}
we have \(4X=YY-(p-q)^2ZZ\), and therefore, since \((p-q)^2=\pm n\), \[4X=YY\mp nZZ\] The upper sign holds when \(n\) is of the form \(4k+1\), the lower when \(n\) is of the form \(4k+3\). This is the theorem whose proof we promised above (Article 124). It will easily be seen that the two highest terms of the function \(Y\) are always \(2x^m+x^{m-1}\), and the highest term of the function \(Z\) is \(x^{m-1}\). The remaining coefficients, which will clearly all be integers, vary with the different nature of the number \(n\) and cannot be given by a general analytic formula.

\textit{Example.} For \(n=17\), the rules of Article 348 give the equation whose roots are the eight roots contained in \((8,1)\) as \[\begin{aligned}x^8-px^7+(4+p+2q)x^6&-(4p+3q)x^5+(6+3p+5q)x^4\\&-(4p+3q)x^3+(4+p+2q)xx-px+1=0\end{aligned}\] whence \[\begin{aligned}R&=x^8+4x^6+6x^4+4xx+1\\S&=-x^7+x^6-4x^5+3x^4-4x^3+xx-x\\T&=2x^6-3x^5+5x^4-3x^3+2xx\end{aligned}\] and hence \[\begin{aligned}Y&=2x^8+x^7+5x^6+7x^5+4x^4\\&\qquad+7x^3+5xx+x+2\\Z&=x^7+x^6+x^5+2x^4+x^3+xx+x\end{aligned}\] Here are a few more examples: {\small\setlength{\arraycolsep}{4pt}\[\begin{array}{c|l|l}n&Y&Z\\\hline 3&2x+1&1\\5&2xx+x+2&x\\7&2x^3+xx-x-2&xx+x.\\11&2x^5+x^4-2x^3+2xx-x-2&x^4+x\\13&2x^6+x^5+4x^4-x^3+4xx+x+2&x^5+x^3+x\\19&2x^9+x^8-4x^7+3x^6+5x^5-5x^4&x^8-x^6+x^5+x^4-x^3+x\\{}&{}-3x^3+4xx-x-2&{}\\23&2x^{11}+x^{10}-5x^9-8x^8-7x^7-4x^6&x^{10}+x^9-x^7-2x^6-2x^5\\{}&{}+4x^5+7x^4+8x^3+5xx-x-2&{}-x^4+xx+x\end{array}\]}
\SourceLeaf{450}

\Needspace{9\baselineskip}\begin{center}\itshape \textit{On the equation for the distribution of the roots \(\Omega\) into three periods.}\end{center}


\Needspace{6\baselineskip}\Article{358}
We proceed to consider the cubic equations that, when \(n\) is of the form \(3k+1\), determine the three sums of \(\frac13(n-1)\) terms constituting the collection \(\Omega\). Let \(g\) be any primitive root modulo \(n\), and put \(\frac13(n-1)=m\); the value thus defined will be an even integer. Then the three sums constituting \(\Omega\) will be \((m,1),\allowbreak (m,g),\allowbreak (m,gg)\), for which we will write \(p,\allowbreak p',\allowbreak p''\), respectively. Clearly, the first contains the roots \([1],\allowbreak [g^3],\allowbreak [g^6]\ldots[g^{n-4}]\), the second \([g],\allowbreak [g^4]\ldots[g^{n-3}]\), and the third \([gg],\allowbreak [g^5]\ldots[g^{n-2}]\). Supposing the required equation to be \[x^3-Axx+Bx-C=0\] we have \[\begin{gathered}A=p+p'+p'',\quad B=pp'+p'p''+pp'',\\C=pp'p''\end{gathered}\]

whence we immediately obtain \(A=-1\). Let the least positive residues of the numbers \(g^3\), \(g^6\ldots g^{n-4}\) modulo \(n\) be, in any order, \(\mathfrak A,\allowbreak \mathfrak B,\allowbreak \mathfrak C\), etc., and let \(\mathfrak K\) be their collection with the number \(1\) added. Similarly, let \(\mathfrak A',\allowbreak \mathfrak B',\allowbreak \mathfrak C'\), etc., be the least residues of the numbers \(g,\allowbreak g^4,\allowbreak g^7\ldots g^{n-3}\), and \(\mathfrak K'\) their collection; finally, let \(\mathfrak A'',\allowbreak \mathfrak B'',\allowbreak \mathfrak C''\), etc., be the least residues of \(gg,\allowbreak g^5,\allowbreak g^8\ldots g^{n-2}\), and \(\mathfrak K''\) their collection. Thus all the numbers in \(\mathfrak K,\allowbreak \mathfrak K',\allowbreak \mathfrak K''\) will be distinct and will coincide with \(1,\allowbreak 2,\allowbreak 3\ldots n-1\). Above all, we must observe here that the number \(n-1\) necessarily occurs in \(\mathfrak K\), since it is easily seen to be the residue of \(g^{\frac{3m}{2}}\). It also follows easily that two numbers such as \(h,\allowbreak n-h\) always occur in the \textit{same} one of the three collections \(\mathfrak K,\allowbreak \mathfrak K',\allowbreak \mathfrak K''\). For if one is the residue of the power \(g^\lambda\), the other will be the residue of the power \(g^{\lambda+\frac{3m}{2}}\), or of \(g^{\lambda-\frac{3m}{2}}\) if \(\lambda>\frac{3m}{2}\). Let \((\mathfrak K\mathfrak K)\) denote the count of numbers in the sequence \(1,\allowbreak 2,\allowbreak 3\ldots n-1\) that, together with the numbers immediately following them, greater by one, are contained in \(\mathfrak K\). Similarly, let \((\mathfrak K\mathfrak K')\) be the count of numbers in the same sequence that are themselves contained in \(\mathfrak K\), while their immediate successors are contained in \(\mathfrak K'\). The meaning of the symbols \((\mathfrak K\mathfrak K''),\allowbreak (\mathfrak K'\mathfrak K),\allowbreak (\mathfrak K'\mathfrak K'),\allowbreak (\mathfrak K'\mathfrak K''),\allowbreak (\mathfrak K''\mathfrak K),\allowbreak (\mathfrak K''\mathfrak K'),\allowbreak (\mathfrak K''\mathfrak K'')\) will then also be clear at once. I now claim, \textit{first}, that \((\mathfrak K\mathfrak K')=(\mathfrak K'\mathfrak K)\). Suppose \(h,\allowbreak h',\allowbreak h''\), etc., are all the numbers of the sequence \(1,\allowbreak 2,\allowbreak 3\ldots n-1\) that are themselves in \(\mathfrak K\), while the next greater numbers \(h+1,\allowbreak h'+1,\allowbreak h''+1\), etc., are in \(\mathfrak K'\), and whose number is therefore \(=(\mathfrak K\mathfrak K')\). It is clear that all the numbers \(n-h-1\), \(n-h'-1,\allowbreak n-h''-1\), etc., are contained in \(\mathfrak K'\), while the next greater numbers \(n-h,\allowbreak n-h'\), etc., are
\SourceLeaf{451}
in \(\mathfrak K\). Since there are altogether \((\mathfrak K'\mathfrak K)\) such numbers, we certainly cannot have \((\mathfrak K'\mathfrak K)<(\mathfrak K\mathfrak K')\). In the same way it is proved that we cannot have \((\mathfrak K\mathfrak K')<(\mathfrak K'\mathfrak K)\), so these numbers must necessarily be equal. Exactly the same reasoning proves \((\mathfrak K\mathfrak K'')=(\mathfrak K''\mathfrak K)\), \((\mathfrak K'\mathfrak K'')=(\mathfrak K''\mathfrak K')\). Second, since every number of \(\mathfrak K\), except the largest, \(n-1\), must have an immediate successor contained either in \(\mathfrak K\), in \(\mathfrak K'\), or in \(\mathfrak K''\), the sum \((\mathfrak K\mathfrak K)+(\mathfrak K\mathfrak K')+(\mathfrak K\mathfrak K'')\) will equal the count of all the numbers in \(\mathfrak K\) diminished by one, namely \(=m-1\). For a similar reason we will have \[\begin{aligned}(\mathfrak K'\mathfrak K)+(\mathfrak K'\mathfrak K')+(\mathfrak K'\mathfrak K'')\\{}=(\mathfrak K''\mathfrak K)+(\mathfrak K''\mathfrak K')+(\mathfrak K''\mathfrak K'')=m\end{aligned}\]

With these preparations, we expand, by the rules of Article 345, the product \(pp'\) into \((m,\mathfrak A'+1)+(m,\mathfrak B'+1)+(m,\mathfrak C'+1)+\text{etc.}\). This expression is easily seen to reduce to \((\mathfrak K'\mathfrak K)p+(\mathfrak K'\mathfrak K')p'+(\mathfrak K'\mathfrak K'')p''\). Since, by Article 345 I, the product \(p'p''\) arises from it by replacing \((m,1),\allowbreak (m,g),\allowbreak (m,gg)\) respectively by \((m,g),\allowbreak (m,gg)\), \((mg^3)\), that is, \(p,\allowbreak p',\allowbreak p''\) respectively by \(p',\allowbreak p'',\allowbreak p\), we obtain \(p'p''=(\mathfrak K'\mathfrak K)p'+(\mathfrak K'\mathfrak K')p''+(\mathfrak K'\mathfrak K'')p\), and in exactly the same way \(p''p=(\mathfrak K'\mathfrak K)p''+(\mathfrak K'\mathfrak K')p+(\mathfrak K'\mathfrak K'')p'\). From this it follows immediately, \textit{first}, that \[B=m(p+p'+p'')=-m\] \textit{second}, since, in the same way as \(pp'\) was expanded before, \(pp''\) also reduces to \((\mathfrak K''\mathfrak K)p+(\mathfrak K''\mathfrak K')p'+(\mathfrak K''\mathfrak K'')p''\), and this expression must be identical with the preceding one, we necessarily have \((\mathfrak K''\mathfrak K)=(\mathfrak K'\mathfrak K')\) and \((\mathfrak K''\mathfrak K'')=(\mathfrak K'\mathfrak K)\). Hence, putting \[\begin{aligned}(\mathfrak K'\mathfrak K'')&=(\mathfrak K''\mathfrak K')=a,\\(\mathfrak K''\mathfrak K'')&=(\mathfrak K'\mathfrak K)=(\mathfrak K\mathfrak K')=b,\\(\mathfrak K'\mathfrak K')&=(\mathfrak K''\mathfrak K)=(\mathfrak K\mathfrak K'')=c\end{aligned}\] we obtain \(m-1=(\mathfrak K\mathfrak K)+(\mathfrak K\mathfrak K')+(\mathfrak K\mathfrak K'')=(\mathfrak K\mathfrak K)+b+c\) and \(a+b+c=m\), whence \((\mathfrak K\mathfrak K)=a-1\). Thus those nine unknown quantities have been reduced to three, \(a,\allowbreak b,\allowbreak c\), or rather, because of the equation \(a+b+c=m\), to two. Finally, it is clear that the square \(pp\) expands into \((m,1+1)+(m,\mathfrak A+1)+(m,\mathfrak B+1)+(m,\mathfrak C+1)+\text{etc.}\). Among the parts of this expression there will be \((m,n)\), which reduces to \((m,0)\), or to \(m\); the remaining parts are easily seen to reduce to \((\mathfrak K\mathfrak K)p+(\mathfrak K\mathfrak K')p'+(\mathfrak K\mathfrak K'')p''\), whence we will have \(pp=m+(a-1)p+bp'+cp''\).

In this way the preceding investigations have given us these four reductions.
\SourceLeaf{452}\[\begin{aligned}pp&=m+(a-1)p+bp'+cp''\\pp'&=bp+cp'+ap''\\pp''&=cp+ap'+bp''\\p'p''&=ap+bp'+cp''\end{aligned}\] Between the three unknowns \(a,\allowbreak b,\allowbreak c\) there is the condition \[a+b+c=m\quad \ldotp\,\ldotp\,\ldotp\,\ldotp\,\ldotp\,\ldotp\,\ldotp\,\ldotp\,\ldotp\,\ldotp\,\ldotp\ (\mathrm{I})\] and it is also certain that they are integers. We deduce \[\begin{aligned}C=p\times p'p''&=app+bpp'+cpp''\\&=am+(aa+bb+cc-a)p\\&\quad+(ab+bc+ac)p'+(ab+bc+ac)p''\end{aligned}\]

But since \(pp'p''\) is a symmetric function of the sums \(p,\allowbreak p',\allowbreak p''\), the coefficients multiplying these in the preceding expression must necessarily be equal (Article 350). This gives the new equation \[aa+bb+cc-a=ab+bc+ac\quad \ldotp\,\ldotp\,\ldotp\ (\mathrm{II})\] and hence \(C=am+(ab+bc+ac)(p+p'+p'')\), or, because of I and \(p+p'+p''=-1\), \[C=aa-bc\quad \ldotp\,\ldotp\,\ldotp\,\ldotp\,\ldotp\,\ldotp\ (\mathrm{III})\]

Now, although \(C\) here depends on three unknowns between which we have only two equations, these equations, together with the condition that \(a,\allowbreak b,\allowbreak c\) are integers, suffice to determine \(C\) completely. To show this, we write equation II as \[\begin{aligned}12a+12b+12c+4&=36aa+36bb+36cc\\&\quad-36ab-36ac-36bc\\&\quad-24a+12b+12c+4\end{aligned}\] By I, the left side is \(=12m+4=4n\); the right side reduces to \[(6a-3b-3c-2)^2+27(b-c)^2\] or, writing \(k\) for \(2a-b-c\), to \((3k-2)^2+27(b-c)^2\). It is thus clear that the number \(4n\) (that is, in general, four times any prime of the form \(3m+1\)) can be represented by the form \(xx+27yy\). This can indeed be deduced without difficulty from the general theory
\SourceLeaf{453}
of binary forms, yet it is quite remarkable that such a decomposition is connected with the values of \(a,\allowbreak b,\allowbreak c\). But the number \(4n\) can always be decomposed in only one way into a square and \(27\) times a square, which we prove as follows\HistoricalNote{*)}{This proposition could be proved more directly from the principles of Section V.}. If we suppose \[4n=tt+27uu=t't'+27u'u'\] we obtain, \textit{first}, \[(tt'-27uu')^2+27(tu'+t'u)^2=16nn\] \textit{second}, \[(tt'+27uu')^2+27(tu'-t'u)^2=16nn\] \textit{third}, \[(tu'+t'u)(tu'-t'u)=4n(u'u'-uu)\] From the third equation it follows that \(n\), since it is a prime number, divides one of the numbers \(tu'+t'u,\allowbreak tu'-t'u\); from the first and second it is clear that each of these numbers is smaller than \(n\). Thus the one divisible by \(n\) must necessarily be \(=0\), and therefore also \(u'u'-uu=0\), whence \(u'u'=uu\) and \(t't'=tt\): the two decompositions do not differ. Thus, if we suppose the decomposition of \(4n\) into a square and \(27\) times a square to be known (it can be found either by the direct method of Section V or by the indirect method given in Articles 323 and 324), namely if \(4n=MM+27NN\), the squares \((3k-2)^2,\allowbreak (b-c)^2\) will be determined, and we will now have obtained two equations in place of equation II. But it is easily seen that not only the square \((3k-2)^2\) but also the root \(3k-2\) itself is fully determined. Since it must necessarily be either \(=+M\) or \(=-M\), the ambiguity is removed by the requirement that \(k\) be an integer. We therefore set \(3k-2=+M\) or \(=-M\), according as \(M\) is of the form \(3z+1\) or \(3z+2\)\HistoricalNote{\(\dagger\))}{Clearly, \(M\) cannot be of the form \(3z\), since otherwise \(4n\) would be divisible by \(3\). — We need not consider here the ambiguity as to whether \(b-c\) should be set \(=N\) or \(=-N\). Nor can it, by the nature of the matter, be removed in any way, since it depends on the choice of the primitive root \(g\): for some primitive roots the difference \(b-c\) is positive, for others negative.}. Now, since \(k=2a-b-c=3a-m\), we have \(a=\frac13(m+k),\allowbreak \ b+c=m-a=\frac13(2m-k)\), whence \[\begin{aligned}C=aa-bc&=aa-\tfrac14(b+c)^2+\tfrac14(b-c)^2\\&=\tfrac19(m+k)^2-\tfrac1{36}(2m-k)^2+\tfrac14NN\\&=\tfrac1{12}kk+\tfrac13km+\tfrac14NN\end{aligned}\] and so all the coefficients of the required equation have been found. \(Q.E.F.\) — This formula
\SourceLeaf{454}
becomes still simpler if for \(NN\) we substitute its value from the equation \((3k-2)^2+27NN=4n=12m+4\), which gives, after calculation, \[C=\tfrac19(m+k+3km)=\tfrac19(m+kn)\] The same value can also be reduced to \((3k-2)NN+k^3-2kk+k-km+m\). Though less suitable for use, this expression immediately shows that \(C\) certainly is an integer, as it should be.

\textit{Example.} For \(n=19\), we have \(4n=49+27\), whence \(3k-2=+7,\allowbreak k=3\), \(C=\frac19(6+57)=7\), and the required equation is \(x^3+xx-6x-7=0\), as above (Article 351). — Similarly, for \(n=7,\allowbreak 13,\allowbreak 31,\allowbreak 37,\allowbreak 43,\allowbreak 61,\allowbreak 67\) the value of \(k\) is found to be, respectively, \(1,\allowbreak -1,\allowbreak 2,\allowbreak -3,\allowbreak -2,\allowbreak 1,\allowbreak -1\), whence \(C=1,\allowbreak -1,\allowbreak 8,\allowbreak -11,\allowbreak -8,\allowbreak 9,\allowbreak -5\).

\clearpage
Although the problem solved in this article is rather intricate, we did not wish to suppress it, both because of the elegance of its solution and because it gave occasion to use various devices that can also be applied with great benefit to other questions\HistoricalNote{*)}{Corollary. Let \(\varepsilon\) be a root of the equation \(x^3-1=0\), and you will have \((p+\varepsilon p'+\varepsilon\varepsilon p'')^3=\frac n2(M+N\surd -27)\). Put \(\dfrac M{\surd 4n}=\cos\varphi,\allowbreak \quad\dfrac{N\surd 27}{\surd 4n}=\sin\varphi\); then \[\begin{gathered}p=-\tfrac13+\tfrac23\cos\tfrac13\varphi\surd n;\\M\equiv+1\ (\operatorname{mod.}3);\quad1\equiv M(1.2.3\ldots m)^3\ (\operatorname{mod.}n)\end{gathered}\] Putting \(3x+1=y\), the equation becomes \(y^3-3ny-Mn=0\).}.

